\section{Antecedentes del proyecto}

\subsection{Situación previa}
\begin{indentedsubsection}
Industria Textil Atlas E.I.R.L. es una empresa fundada en el año 2004 en la ciudad de Arequipa por Víctor Huayllaro, operando comercialmente desde la céntrica zona de "San Camilo". A lo largo de más de 20 años de trayectoria, la empresa ha consolidado una marca local respetable basada en la calidad y durabilidad de sus productos de protección deportiva y prendas de licra. A pesar de llevar la palabra "Textil" en su denominación social, la empresa no realiza confección ni costura directa en sus instalaciones; en su lugar, terceriza por completo la fabricación de sus productos a talleres externos de confección. Su modelo de negocio se basa en la distribución y comercialización mayorista y minorista de una amplia gama de artículos de soporte, protección y accesorios deportivos (rodilleras, musleras, tobilleras, muñequeras, coderas, pantorrilleras, menisqueras, soportes de rodilla, soportes y moldeadores de cintura, suspensores, pelotas de fútbol, balones de balonmano, redes de básquet, sets de ping pong, mitones, mangas de vóley, cintas kinesiológicas, ankles, polos y pantalonetas de licra, y bandas de resistencia).

Hasta antes del inicio del proyecto, la gestión del negocio se caracterizaba por un modelo tradicional y empírico:

\begin{itemize}
    \item \textbf{Gestión de ventas informal:} la tienda física de San Camilo y el canal digital de WhatsApp operaban de forma desarticulada. El personal de ventas registraba las transacciones de forma manual en cuadernos o mediante anotaciones sueltas ("papelitos"), sin una interfaz rápida para procesar las ventas en el mostrador.
    \item \textbf{Cálculo de márgenes y fijación de precios analógica:} la gerencia carecía de una herramienta dinámica para monitorear las variaciones en los costos de adquisición de los proveedores externos. Esto dificultaba la fijación de precios de venta finales y el cálculo exacto de la rentabilidad deseada, obligando a manejar la estructura de precios de memoria o en registros manuales dispersos.
    \item \textbf{Verificación visual de existencias:} el control de stock en almacén se realizaba mediante revisión visual directa de los estantes y anotaciones físicas en libretas, dependiendo críticamente de la memoria y presencia del personal de almacén para determinar si había unidades disponibles de un determinado artículo o talla.
\end{itemize}
\end{indentedsubsection}

\subsection{Problemática identificada}
\begin{indentedsubsection}
El diagnóstico de los procesos de comercialización y almacenamiento reveló brechas críticas que afectaban la rentabilidad, la agilidad en la atención y la eficiencia del capital de trabajo de la empresa:

\begin{itemize}
    \item \textbf{A. Descoordinación de Ventas vs. Almacén:} debido a la falta de un sistema de control de existencias en tiempo real, el personal de ventas en tienda o WhatsApp confirmaba con frecuencia transacciones de productos agotados físicamente. Esto provocaba llamadas de rectificación a los clientes, cancelaciones y retrasos que deterioraban la experiencia de compra y la lealtad hacia la marca.
    \item \textbf{B. Lentitud extrema en la atención al cliente:} el proceso de búsqueda, validación y despacho de artículos en la tienda física de San Camilo tomaba entre 15 y 20 minutos por cliente (frente a un óptimo de 5 minutos). Este retraso se debía a la ausencia de una codificación estandarizada (SKU) y al desorden físico en las estanterías, obligando a los vendedores a buscar manualmente talla por talla y artículo por artículo.
    \item \textbf{C. Descontrol financiero y operativo de inventarios:} la falta de automatización mantenía un margen de error superior al 30\% en los registros de almacén. Esto ocasionaba compras reactivas a los talleres externos, sobreabastecimiento de artículos de baja rotación (congelando capital de trabajo) y pérdidas por productos que quedaban obsoletos o descatalogados. De manera simultánea, ocurrían quiebres de stock en accesorios deportivos clave.
    \item \textbf{D. Dificultad en la gestión de precios y márgenes por Gerencia:} al no centralizarse el registro de costos de compra de los proveedores en un sistema integrado, la gerencia no podía ajustar ágilmente los precios de venta para asegurar márgenes de ganancia estables frente a las variaciones de los insumos y costos de tercerización, centralizando excesivamente la toma de decisiones financieras en estimaciones del gerente.
\end{itemize}
\end{indentedsubsection}
